% Generated draft: requires mathematical and source review.
\documentclass[11pt]{article}
\usepackage[margin=25mm]{geometry}
\usepackage{fontspec}
\setmainfont{Latin Modern Roman}
\usepackage{amsmath,amssymb,mathtools,unicode-math}
\setmathfont{Latin Modern Math}
\usepackage{enumitem,needspace,xurl,hyperref}
\setlength{\parindent}{0pt}
\setlength{\parskip}{5pt}
\setlength{\emergencystretch}{3em}
\newcommand{\N}{\mathbb{N}}
\newcommand{\Z}{\mathbb{Z}}
\newcommand{\Q}{\mathbb{Q}}
\newcommand{\R}{\mathbb{R}}
\newcommand{\C}{\mathbb{C}}
\newcommand{\F}{\mathbb{F}}
\newcommand{\A}{\mathbb{A}}
\newcommand{\PP}{\mathbb{P}}
\newcommand{\E}{\mathbb{E}}
\newcommand{\eps}{\varepsilon}
\newcommand{\dd}{\,\mathrm{d}}
\begin{document}
\section*{UnsolvedMath ID: 30004717}
\textbf{Source problem code:} catalogue entry
\subsection{Definitions and mathematical statement}
Let \(n\geq 1\) be an integer, and let \(Y\) be a closed smooth manifold of dimension \(n-1\), with \(n-1\neq 4\). Here closed means compact and without boundary. Assume that \(Y\) admits no smooth Riemannian metric \(h\) with strictly positive scalar curvature everywhere. Set
\[
M=Y\times[-1,1],\qquad \partial_-M=Y\times\{-1\},\qquad \partial_+M=Y\times\{1\}.
\]
For a smooth Riemannian metric \(g\) on \(M\), smooth up to its boundary, define its band width by
\[
\operatorname{width}(M,g)=\operatorname{dist}_g(\partial_-M,\partial_+M)
=\inf_{\gamma} L_g(\gamma),
\]
where the infimum is over piecewise smooth curves \(\gamma:[0,1]\to M\) with \(\gamma(0)\in\partial_-M\) and \(\gamma(1)\in\partial_+M\), and \(L_g\) denotes Riemannian length.

The supplied conjecture asserts that, for every such \(Y\), every real number \(\sigma>0\), and every such metric \(g\) satisfying
\[
\operatorname{Scal}_g(p)\geq\sigma\qquad\text{for all }p\in M,
\]
one has
\[
\operatorname{width}(M,g)\leq 2\pi\sqrt{\frac{n-1}{\sigma n}}.
\]
The source states no connectedness, orientability, spin, or boundary mean-curvature hypothesis; none is added here. It gives no further lower dimension bound beyond that implicit in \(\dim Y=n-1\). Its supplied literature assessment concerns partial results and does not establish this full conjecture. In particular, the phrase ``orientable dimensions through seven'' in that assessment does not specify whether it refers to \(\dim Y\) or \(\dim M\), nor does it give the precise additional hypotheses of those results.
\subsection{Short English statement}
The conjecture bounds the distance between the two boundary components of \(Y\times[-1,1]\) whenever its scalar curvature is at least \(\sigma>0\) and \(Y\), of dimension other than four, admits no positive-scalar-curvature metric. The proposed bound is \(2\pi\sqrt{(n-1)/(\sigma n)}\). The supplied record labels the problem partially solved; this draft requires human mathematical review.
\subsection{Sources}
\begin{itemize}
\item Supplied local UnsolvedMath registry record, permanent ID \(30004717\), title: \emph{Sharp Width Bounds for Positive-Scalar-Curvature Bands}; collection: Oberwolfach Reports Open Problems; exact source problem code: \texttt{catalogue entry}. The record gives proposed year 2021 and names no individual proposer.
\item Supplied report citation: \emph{Analysis, Geometry and Topology of Positive Scalar Curvature Metrics} (2022), Oberwolfach Reports, pp.\ 1579--1649, DOI \texttt{10.4171/owr/2021/30}. \url{https://doi.org/10.4171/owr/2021/30}. The background describes these reports as documenting mathematical workshops and their open-problem discussions.
\item Supplied literature reference: Cecchini, Rade and Zeidler, \emph{Journal of Topology}. \url{https://londmathsoc.onlinelibrary.wiley.com/doi/full/10.1112/topo.12303}. The record supplies no article title or theorem number.
\item Supplied literature reference: Rade, doctoral thesis on band width. \url{https://opus.bibliothek.uni-augsburg.de/opus4/frontdoor/index/index/year/2023/docId/99573}. The record supplies no full thesis title or theorem number.
\item Literature assessment dated exactly \texttt{2026-08-22}: supplied dated curation derived from the OWR record, not independently verified here. It labels the problem partially solved and reports that the band-width inequality is established in important spin and index-theoretic cases, including orientable dimensions through seven, while no full all-dimension and all-topology solution is known. The supplied background explicitly requests independent verification of time-sensitive status and cited claims.
\item The registry labels the statement \texttt{corrected\_verified} and says that the band-width bound and its dimension and curvature hypotheses were normalized without changing the conjecture. These are supplied curation claims, not independent mathematical verification by this draft.
\end{itemize}
\end{document}
